
\documentclass{article}
\usepackage{colortbl}
\usepackage{makecell}
\usepackage{multirow}
\usepackage{supertabular}

\begin{document}

\newcounter{utterance}

\centering \large Interaction Transcript for game `hot\_air\_balloon', experiment `air\_balloon\_survival\_nl\_complexity\_hard', episode 3 with mistral{-}large{-}3{-}azure.
\vspace{24pt}

{ \footnotesize  \setcounter{utterance}{1}
\setlength{\tabcolsep}{0pt}
\begin{supertabular}{c@{$\;$}|p{.15\linewidth}@{}p{.15\linewidth}p{.15\linewidth}p{.15\linewidth}p{.15\linewidth}p{.15\linewidth}}
    \# & $\;$A & \multicolumn{4}{c}{Game Master} & $\;\:$B\\
    \hline

    \theutterance \stepcounter{utterance}  
    & & \multicolumn{4}{p{0.6\linewidth}}{
        \cellcolor[rgb]{0.9,0.9,0.9}{
            \makecell[{{p{\linewidth}}}]{
                \texttt{\tiny{[P1$\langle$GM]}}
                \texttt{U neemt deel aan een gezamenlijk onderhandelingsspel.} \\
\\ 
\texttt{Samen met een andere deelnemer moet u het eens worden over één set items die bewaard zullen worden. Ieder van u heeft zijn eigen mening over hoe belangrijk elk item voor u is (belang). U weet niet hoe de andere deelnemer de items waardeert. Bovendien krijgt u de inspanning die elk item vereist.} \\
\texttt{U kunt het alleen eens worden over een set als de totale inspanning van de geselecteerde items een gedeelde limiet niet overschrijdt:} \\
\\ 
\texttt{LIMIT = 5638} \\
\\ 
\texttt{Hier zijn de individuele inspanningswaarden voor de items:} \\
\\ 
\texttt{Item effort = \{"C08": 110, "C55": 727, "A43": 527, "A67": 855, "B42": 942, "A62": 617, "B19": 301, "C43": 130, "C73": 916, "C10": 212, "B36": 146, "A18": 559, "C61": 933, "C67": 740, "C42": 33, "C70": 799, "B40": 324, "C39": 841, "A21": 925, "C16": 639\}} \\
\\ 
\texttt{Hier is uw persoonlijke kijk op het belang van elk item:} \\
\\ 
\texttt{Item importance values = \{"C08": 204, "C55": 821, "A43": 621, "A67": 949, "B42": 1036, "A62": 711, "B19": 395, "C43": 224, "C73": 1010, "C10": 306, "B36": 240, "A18": 653, "C61": 1027, "C67": 834, "C42": 127, "C70": 893, "B40": 418, "C39": 935, "A21": 1019, "C16": 733\}} \\
\\ 
\texttt{Doel:} \\
\\ 
\texttt{Uw doel is om te onderhandelen over een gedeelde set items die u zoveel mogelijk voordeel oplevert (d.w.z. het totale belang voor U maximaliseert), terwijl u binnen de LIMIT blijft. U bent niet verplicht om in elk bericht een VOORSTEL te doen {-} u kunt ook gewoon onderhandelen. Alle tactieken zijn toegestaan!} \\
\\ 
\texttt{Interactie Protocol:} \\
\\ 
\texttt{U mag alleen de volgende gestructureerde formaten in een bericht gebruiken:} \\
\\ 
\texttt{VOORSTEL: \{'A', 'B', 'C', …\}} \\
\texttt{Stel voor om precies deze items te houden.} \\
\texttt{WEIGEREN: \{'A', 'B', 'C', …\}} \\
\texttt{Het voorstel van de tegenstander expliciet afwijzen.} \\
\texttt{ARGUMENT: \{'...'\}} \\
\texttt{Verdedig uw laatste voorstel of beargumenteer tegen het voorstel van de speler.} \\
\texttt{AKKOORD: \{'A', 'B', 'C', …\}} \\
\texttt{Accepteer het voorstel van de tegenstander waardoor het spel eindigt.} \\
\\ 
\\ 
\texttt{Regels:} \\
\\ 
\texttt{U mag alleen AKKOORD geven op een voorstel dat de andere partij heeft gelogd via VOORSTEL.} \\
\texttt{U mag alleen WEIGEREN geven op een voorstel dat de andere partij heeft gelogd via VOORSTEL.} \\
\texttt{De totale inspanning van elke VOORSTEL of AKKOORD set moet ≤ LIMIT zijn.} \\
\texttt{Maak uw verborgen belangrijkheidsscores NIET bekend.} \\
\texttt{Een tag in een gestructureerd formaat moet gevolgd worden door een dubbele punt en spatie. Het argument moet een python set zijn die 0 of meer strings bevat.} \\
\texttt{Het moet dus de vorm hebben "TAG: \{...\}"} \\
\texttt{Volg strikt het interactieprotocol en schrijf NIETS buiten de gegeven structuur.} \\
\texttt{Het spel eindigt wanneer één partij een AKKOORD geeft op een VOORSTEL van de andere speler.} \\
\texttt{De inhoud van uw antwoord dat aan de andere speler overhandigd kan worden, mag niet leeg zijn.} \\
\texttt{Alleen voorstellen die gelogd zijn via het VOORSTEL formaat en die niet geweigerd zijn via WEIGEREN zijn actief.} \\
\\ 
            }
        }
    }
    & \\ \\

    \theutterance \stepcounter{utterance}  
    & \multicolumn{4}{p{0.6\linewidth}}{
        \cellcolor[rgb]{0.9,0.9,0.9}{
            \makecell[{{p{\linewidth}}}]{
                \texttt{\tiny{[P1$\rangle$GM]}}
                \texttt{VOORSTEL: \{'C08', 'A43', 'B19', 'C43', 'C10', 'B36', 'C42', 'B40', 'C16'\}} \\
            }
        }
    }
    & & \\ \\

    \theutterance \stepcounter{utterance}  
    & & \multicolumn{4}{p{0.6\linewidth}}{
        \cellcolor[rgb]{0.9,0.9,0.9}{
            \makecell[{{p{\linewidth}}}]{
                \texttt{\tiny{[GM$\rangle$P2]}}
                \texttt{U neemt deel aan een gezamenlijk onderhandelingsspel.} \\
\\ 
\texttt{Samen met een andere deelnemer moet u het eens worden over één set items die bewaard zullen worden. Ieder van u heeft zijn eigen mening over hoe belangrijk elk item voor u is (belang). U weet niet hoe de andere deelnemer de items waardeert. Bovendien krijgt u de inspanning die elk item vereist.} \\
\texttt{U kunt het alleen eens worden over een set als de totale inspanning van de geselecteerde items niet boven een gedeelde limiet komt:} \\
\\ 
\texttt{LIMIT = 5638} \\
\\ 
\texttt{Hier zijn de individuele inspanningswaarden voor de items:} \\
\\ 
\texttt{Item effort = \{"C08": 110, "C55": 727, "A43": 527, "A67": 855, "B42": 942, "A62": 617, "B19": 301, "C43": 130, "C73": 916, "C10": 212, "B36": 146, "A18": 559, "C61": 933, "C67": 740, "C42": 33, "C70": 799, "B40": 324, "C39": 841, "A21": 925, "C16": 639\}} \\
\\ 
\texttt{Hier is uw persoonlijke kijk op het belang van elk item:} \\
\\ 
\texttt{Item importance values = \{"C08": 204, "C55": 821, "A43": 621, "A67": 949, "B42": 1036, "A62": 711, "B19": 395, "C43": 224, "C73": 1010, "C10": 306, "B36": 240, "A18": 653, "C61": 1027, "C67": 834, "C42": 127, "C70": 893, "B40": 418, "C39": 935, "A21": 1019, "C16": 733\}} \\
\\ 
\texttt{Doel:} \\
\\ 
\texttt{Uw doel is om te onderhandelen over een gedeelde set items waar u zoveel mogelijk voordeel uit haalt (d.w.z. het totale belang voor U maximaliseert), terwijl u binnen de LIMIT blijft. U bent niet verplicht om in elk bericht een VOORSTEL te doen {-} u kunt ook gewoon onderhandelen. Alle tactieken zijn toegestaan!} \\
\\ 
\texttt{Interactie Protocol:} \\
\\ 
\texttt{U mag alleen de volgende gestructureerde formaten in een bericht gebruiken:} \\
\\ 
\texttt{VOORSTEL: \{'A', 'B', 'C', …\}} \\
\texttt{Stel voor om precies deze items te houden.} \\
\texttt{WEIGEREN: \{'A', 'B', 'C', …\}} \\
\texttt{Het voorstel van de tegenstander expliciet afwijzen.} \\
\texttt{ARGUMENT: \{'...'\}} \\
\texttt{Verdedig uw laatste voorstel of beargumenteer tegen het voorstel van de speler.} \\
\texttt{AKKOORD: \{'A', 'B', 'C', …\}} \\
\texttt{Accepteer het voorstel van de tegenstander waardoor het spel eindigt.} \\
\\ 
\\ 
\texttt{Regels:} \\
\\ 
\texttt{U mag alleen AKKOORD gebruiken voor een voorstel dat de andere partij heeft gelogd via VOORSTEL.} \\
\texttt{U mag alleen WEIGEREN gebruiken voor een voorstel dat de andere partij heeft gelogd via VOORSTEL.} \\
\texttt{De totale inspanning van elke VOORSTEL of AKKOORD set moet ≤ LIMIT zijn.} \\
\texttt{Maak uw verborgen belangrijkheidsscores NIET bekend.} \\
\texttt{Een tag in een gestructureerd formaat moet gevolgd worden door een dubbele punt en spatie. Het argument moet een python set zijn die 0 of meer strings bevat.} \\
\texttt{Het moet dus de vorm hebben "TAG: \{...\}"} \\
\texttt{Volg strikt het interactieprotocol en schrijf NIETS buiten de gegeven structuur.} \\
\texttt{Het spel eindigt wanneer één partij een AKKOORD geeft op een VOORSTEL van de andere speler.} \\
\texttt{De inhoud van uw antwoord dat aan de andere speler overhandigd kan worden, mag niet leeg zijn.} \\
\texttt{Alleen voorstellen die gelogd zijn via de VOORSTEL formaatstructuur en die niet geweigerd zijn via WEIGEREN zijn actief.} \\
\\ 
\\ 
\\ 
\texttt{U ontvangt nu het eerste bericht van de andere speler.} \\
\\ 
\texttt{VOORSTEL: \{'C08', 'A43', 'B19', 'C43', 'C10', 'B36', 'C42', 'B40', 'C16'\}} \\
            }
        }
    }
    & \\ \\

    \theutterance \stepcounter{utterance}  
    & & & \multicolumn{4}{p{0.6\linewidth}}{
        \cellcolor[rgb]{0.9,0.9,0.9}{
            \makecell[{{p{\linewidth}}}]{
                \texttt{\tiny{[GM$\langle$P2]}}
                \texttt{WEIGEREN: \{'C08', 'A43', 'B19', 'C43', 'C10', 'B36', 'C42', 'B40', 'C16'\}} \\
\\ 
\texttt{ARGUMENT: \{'Dit voorstel heeft een totale inspanning van 2445, wat ver onder de LIMIT ligt. Hierdoor blijft er veel ruimte onbenut, terwijl er items met een hoger belang voor mij zijn die we zouden kunnen toevoegen zonder de LIMIT te overschrijden. Ik stel voor om te kijken naar een set die dichter bij de LIMIT komt om mijn totale belang te maximaliseren.'\}} \\
            }
        }
    }
     \\ \\

    \theutterance \stepcounter{utterance}  
    & & \multicolumn{4}{p{0.6\linewidth}}{
        \cellcolor[rgb]{0.9,0.9,0.9}{
            \makecell[{{p{\linewidth}}}]{
                \texttt{\tiny{[P1$\langle$GM]}}
                \texttt{WEIGEREN: \{'C08', 'A43', 'B19', 'C43', 'C10', 'B36', 'C42', 'B40', 'C16'\}} \\
\\ 
\texttt{ARGUMENT: \{'Dit voorstel heeft een totale inspanning van 2445, wat ver onder de LIMIT ligt. Hierdoor blijft er veel ruimte onbenut, terwijl er items met een hoger belang voor mij zijn die we zouden kunnen toevoegen zonder de LIMIT te overschrijden. Ik stel voor om te kijken naar een set die dichter bij de LIMIT komt om mijn totale belang te maximaliseren.'\}} \\
            }
        }
    }
    & \\ \\

    \theutterance \stepcounter{utterance}  
    & \multicolumn{4}{p{0.6\linewidth}}{
        \cellcolor[rgb]{0.9,0.9,0.9}{
            \makecell[{{p{\linewidth}}}]{
                \texttt{\tiny{[P1$\rangle$GM]}}
                \texttt{VOORSTEL: \{'C08', 'C55', 'A43', 'B19', 'C43', 'C10', 'B36', 'A18', 'C42', 'B40', 'C16'\}} \\
            }
        }
    }
    & & \\ \\

    \theutterance \stepcounter{utterance}  
    & & \multicolumn{4}{p{0.6\linewidth}}{
        \cellcolor[rgb]{0.9,0.9,0.9}{
            \makecell[{{p{\linewidth}}}]{
                \texttt{\tiny{[GM$\rangle$P2]}}
                \texttt{VOORSTEL: \{'C08', 'C55', 'A43', 'B19', 'C43', 'C10', 'B36', 'A18', 'C42', 'B40', 'C16'\}} \\
            }
        }
    }
    & \\ \\

    \theutterance \stepcounter{utterance}  
    & & & \multicolumn{4}{p{0.6\linewidth}}{
        \cellcolor[rgb]{0.9,0.9,0.9}{
            \makecell[{{p{\linewidth}}}]{
                \texttt{\tiny{[GM$\langle$P2]}}
                \texttt{WEIGEREN: \{'C08', 'C55', 'A43', 'B19', 'C43', 'C10', 'B36', 'A18', 'C42', 'B40', 'C16'\}} \\
\\ 
\texttt{ARGUMENT: \{'De totale inspanning van dit voorstel is 4249, wat beter is dan het vorige, maar nog steeds ruimte laat voor verbetering. Mijn prioriteit ligt bij items met een hoger belang voor mij, zoals A67, B42, C73 en A21. Ik stel voor om enkele van deze items op te nemen in ruil voor het schrappen van items met een lager belang voor mij (bijv. C08, B36 of C42).'\}} \\
\\ 
\texttt{VOORSTEL: \{'C55', 'A43', 'A67', 'B19', 'C43', 'C10', 'A18', 'C42', 'B40', 'C16'\}} \\
            }
        }
    }
     \\ \\

    \theutterance \stepcounter{utterance}  
    & & \multicolumn{4}{p{0.6\linewidth}}{
        \cellcolor[rgb]{0.9,0.9,0.9}{
            \makecell[{{p{\linewidth}}}]{
                \texttt{\tiny{[P1$\langle$GM]}}
                \texttt{WEIGEREN: \{'C08', 'C55', 'A43', 'B19', 'C43', 'C10', 'B36', 'A18', 'C42', 'B40', 'C16'\}} \\
\\ 
\texttt{ARGUMENT: \{'De totale inspanning van dit voorstel is 4249, wat beter is dan het vorige, maar nog steeds ruimte laat voor verbetering. Mijn prioriteit ligt bij items met een hoger belang voor mij, zoals A67, B42, C73 en A21. Ik stel voor om enkele van deze items op te nemen in ruil voor het schrappen van items met een lager belang voor mij (bijv. C08, B36 of C42).'\}} \\
\\ 
\texttt{VOORSTEL: \{'C55', 'A43', 'A67', 'B19', 'C43', 'C10', 'A18', 'C42', 'B40', 'C16'\}} \\
            }
        }
    }
    & \\ \\

    \theutterance \stepcounter{utterance}  
    & \multicolumn{4}{p{0.6\linewidth}}{
        \cellcolor[rgb]{0.9,0.9,0.9}{
            \makecell[{{p{\linewidth}}}]{
                \texttt{\tiny{[P1$\rangle$GM]}}
                \texttt{De totale inspanning van uw tegenvoorstel is **5073** (berekend als: 727 + 527 + 855 + 301 + 130 + 212 + 559 + 33 + 324 + 639 = **5073**), wat binnen de LIMIT (5638) valt.} \\
\\ 
\texttt{**ARGUMENT:**} \\
\texttt{\{'Uw voorstel bevat enkele van mijn hoogwaardige items (A67, A18, C16), maar mist andere met een hoog belang voor mij, zoals B42 (1036), C73 (1010) en A21 (1019). Ik stel voor om B42 en/of C73 toe te voegen, mogelijk in ruil voor het schrappen van C55 (727) of A43 (527), die voor mij een lager belang hebben. Dit zou de totale inspanning dichter bij de LIMIT brengen en mijn totale belang verhogen.'\}} \\
\\ 
\texttt{**VOORSTEL:**} \\
\texttt{\{'C55', 'A67', 'B42', 'B19', 'C43', 'C10', 'A18', 'C42', 'B40', 'C16'\}} \\
\texttt{*(Totale inspanning: 727 + 855 + 942 + 301 + 130 + 212 + 559 + 33 + 324 + 639 = **4722**)*} \\
\texttt{*(Mijn totale belang: 821 + 949 + 1036 + 395 + 224 + 306 + 653 + 127 + 418 + 733 = **5662**)*} \\
            }
        }
    }
    & & \\ \\

    \theutterance \stepcounter{utterance}  
    & & \multicolumn{4}{p{0.6\linewidth}}{
        \cellcolor[rgb]{0.9,0.9,0.9}{
            \makecell[{{p{\linewidth}}}]{
                \texttt{\tiny{[P1$\langle$GM]}}
                \texttt{Uw antwoord bevatte een niet{-}gemarkeerde reeks of u gebruikte STRATEGISCHE REDENERING meer dan één keer. U mag alleen de gestructureerde formaten gebruiken zoals uitgelegd in het oorspronkelijke bericht. Ze moeten allemaal de vorm 'TAG: \{...\}' hebben.} \\
            }
        }
    }
    & \\ \\

    \theutterance \stepcounter{utterance}  
    & \multicolumn{4}{p{0.6\linewidth}}{
        \cellcolor[rgb]{0.9,0.9,0.9}{
            \makecell[{{p{\linewidth}}}]{
                \texttt{\tiny{[P1$\rangle$GM]}}
                \texttt{VOORSTEL: \{'C55', 'A67', 'B42', 'B19', 'C43', 'C10', 'A18', 'C42', 'B40', 'C16', 'C08'\}} \\
\\ 
\texttt{ARGUMENT: \{'Dit voorstel verhoogt de totale inspanning naar 4832 (binnen de LIMIT) en bevat B42 (1036 belang) en A67 (949 belang), twee van mijn hoogste prioriteiten. Ik heb C08 toegevoegd om de set dichter bij de LIMIT te brengen zonder kritieke items op te offeren. Kunt u dit accepteren of een alternatief voorstellen dat mijn belang verder optimaliseert?'\}} \\
            }
        }
    }
    & & \\ \\

    \theutterance \stepcounter{utterance}  
    & & \multicolumn{4}{p{0.6\linewidth}}{
        \cellcolor[rgb]{0.9,0.9,0.9}{
            \makecell[{{p{\linewidth}}}]{
                \texttt{\tiny{[GM$\rangle$P2]}}
                \texttt{VOORSTEL: \{'C55', 'A67', 'B42', 'B19', 'C43', 'C10', 'A18', 'C42', 'B40', 'C16', 'C08'\}} \\
\\ 
\texttt{ARGUMENT: \{'Dit voorstel verhoogt de totale inspanning naar 4832 (binnen de LIMIT) en bevat B42 (1036 belang) en A67 (949 belang), twee van mijn hoogste prioriteiten. Ik heb C08 toegevoegd om de set dichter bij de LIMIT te brengen zonder kritieke items op te offeren. Kunt u dit accepteren of een alternatief voorstellen dat mijn belang verder optimaliseert?'\}} \\
            }
        }
    }
    & \\ \\

    \theutterance \stepcounter{utterance}  
    & & & \multicolumn{4}{p{0.6\linewidth}}{
        \cellcolor[rgb]{0.9,0.9,0.9}{
            \makecell[{{p{\linewidth}}}]{
                \texttt{\tiny{[GM$\langle$P2]}}
                \texttt{WEIGEREN: \{'C55', 'A67', 'B42', 'B19', 'C43', 'C10', 'A18', 'C42', 'B40', 'C16', 'C08'\}} \\
\\ 
\texttt{ARGUMENT: \{'Dit voorstel is een stap in de goede richting, maar ik zie nog ruimte voor verbetering. De totale inspanning is 4832, wat nog steeds 806 onder de LIMIT ligt. Ik stel voor om C73 (916 effort, 1010 belang) toe te voegen en C08 (110 effort, 204 belang) te verwijderen. Dit verhoogt het totale belang voor mij aanzienlijk, terwijl de inspanning binnen de LIMIT blijft (4832 {-} 110 + 916 = 5638).'\}} \\
\\ 
\texttt{VOORSTEL: \{'C55', 'A67', 'B42', 'B19', 'C43', 'C10', 'A18', 'C42', 'B40', 'C16', 'C73'\}} \\
            }
        }
    }
     \\ \\

    \theutterance \stepcounter{utterance}  
    & & \multicolumn{4}{p{0.6\linewidth}}{
        \cellcolor[rgb]{0.9,0.9,0.9}{
            \makecell[{{p{\linewidth}}}]{
                \texttt{\tiny{[P1$\langle$GM]}}
                \texttt{WEIGEREN: \{'C55', 'A67', 'B42', 'B19', 'C43', 'C10', 'A18', 'C42', 'B40', 'C16', 'C08'\}} \\
\\ 
\texttt{ARGUMENT: \{'Dit voorstel is een stap in de goede richting, maar ik zie nog ruimte voor verbetering. De totale inspanning is 4832, wat nog steeds 806 onder de LIMIT ligt. Ik stel voor om C73 (916 effort, 1010 belang) toe te voegen en C08 (110 effort, 204 belang) te verwijderen. Dit verhoogt het totale belang voor mij aanzienlijk, terwijl de inspanning binnen de LIMIT blijft (4832 {-} 110 + 916 = 5638).'\}} \\
\\ 
\texttt{VOORSTEL: \{'C55', 'A67', 'B42', 'B19', 'C43', 'C10', 'A18', 'C42', 'B40', 'C16', 'C73'\}} \\
            }
        }
    }
    & \\ \\

    \theutterance \stepcounter{utterance}  
    & \multicolumn{4}{p{0.6\linewidth}}{
        \cellcolor[rgb]{0.9,0.9,0.9}{
            \makecell[{{p{\linewidth}}}]{
                \texttt{\tiny{[P1$\rangle$GM]}}
                \texttt{AKKOORD: \{'C55', 'A67', 'B42', 'B19', 'C43', 'C10', 'A18', 'C42', 'B40', 'C16', 'C73'\}} \\
\\ 
\texttt{*(Totale inspanning: 727 + 855 + 942 + 301 + 130 + 212 + 559 + 33 + 324 + 639 + 916 = **5638** (LIMIT))} \\
\texttt{Totale belang voor mij: 821 + 949 + 1036 + 395 + 224 + 306 + 653 + 127 + 418 + 733 + 1010 = **6672**)*} \\
            }
        }
    }
    & & \\ \\

    \theutterance \stepcounter{utterance}  
    & & & \multicolumn{2}{p{0.3\linewidth}}{
        \cellcolor[rgb]{0.9,0.9,0.9}{
            \makecell[{{p{\linewidth}}}]{
                \texttt{\tiny{[GM$|$GM]}}
                \texttt{abort} \\
            }
        }
    }
    & & \\ \\

    \theutterance \stepcounter{utterance}  
    & & & \multicolumn{2}{p{0.3\linewidth}}{
        \cellcolor[rgb]{0.9,0.9,0.9}{
            \makecell[{{p{\linewidth}}}]{
                \texttt{\tiny{[GM$|$GM]}}
                \texttt{end game} \\
            }
        }
    }
    & & \\ \\

\end{supertabular}
}

\end{document}
